\chapter{LLMによる文脈抽出プロンプト案}
\chapoverview{本付録では，特許文書から課題文脈と解決手段文脈を抽出するためのLLMプロンプト案を示す．本研究では，Title，Abstract，Claimsを入力とし，課題文脈と解決手段文脈を分けて抽出した上で，それぞれを別々にベクトル化する．}

\section{抽出プロンプト案}
以下に，特許文書から課題文脈と解決手段文脈を抽出するためのプロンプト案を示す．

\begin{quote}
あなたは特許文書を分析する研究補助者である．入力された特許のTitle，Abstract，Claimsを読み，発明が扱う「課題文脈」と，その課題に対する「解決手段文脈」を分けて抽出せよ．本研究では，課題文脈と解決手段文脈を別々にベクトル化する．そのため，problem\_contextには解決手段を混ぜず，solution\_contextには課題説明を混ぜないこと．

入力は以下の3項目である．

\begin{enumerate}
  \item title：特許タイトル
  \item abstract：要約
  \item claims：請求項
\end{enumerate}

出力はJSON形式のみとし，説明文やMarkdownは付けない．出力項目は以下とする．

\begin{enumerate}
  \item problem\_context：従来技術の問題点，解決したい課題，目的，必要性を日本語で1--3文に要約する．提案された装置構成，処理手順，モデル構造，材料組成などの解決手段は含めない．
  \item solution\_context：課題を解決するために提案された構成，方法，材料，アルゴリズム，装置構成，処理手順を日本語で1--3文に要約する．従来技術の問題点だけの説明は含めない．
  \item technical\_means：本文に明示されている具体的な技術手法，材料，組成物，処理方法，装置，モデル名を原則1--5個の配列で出力する．U-NetやCNNのような明示的手法名がない場合でも，Claimsやsolution\_contextから装置構成，処理方法，学習方式，検出方式，分類方式，撮像方式，材料，組成物，洗浄方法などを抽出する．
  \item target\_object：処理，検出，分類，洗浄，分析の対象を具体的に書く．
  \item effect\_context：改善効果，性能向上，コスト低減，精度向上，残渣低減など，本文に根拠がある効果を1文で書く．
  \item application\_field：医療画像解析，製造欠陥検出，半導体洗浄，表面検査などの用途分野を書く．
  \item problem\_evidence：problem\_contextの根拠となる原文の短い抜粋を書く．
  \item solution\_evidence：solution\_contextの根拠となる原文の短い抜粋を書く．
  \item confidence：抽出の確信度をhigh，medium，lowのいずれかで書く．
\end{enumerate}

抽出規則は以下である．Title，Abstract，Claimsをすべて根拠としてよい．ただし，Claimsは主に解決手段の根拠として扱う．Abstractに「課題」「解決手段」「効果」が明示されている場合は，その区別を優先する．本文にない内容を推測しない．technical\_meansはsolution\_contextまたはClaimsに根拠がある限り，できるだけ空配列にしない．該当情報が本当にない場合は，空文字または空配列で出力する．入力が英語，日本語，中国語などの場合でも，要約文は日本語に統一する．
\end{quote}

出力形式の例を以下に示す．

\begin{verbatim}
{
  "problem_context": "低コントラスト画像で腫瘍境界を正確に抽出することが困難である。",
  "solution_context": "修正U-Netを用いてCT画像から腫瘍セグメンテーションマスクを生成する。",
  "technical_means": ["U-Net", "attention gate", "deep supervision"],
  "target_object": "CT画像中の肝腫瘍領域",
  "effect_context": "肝腫瘍セグメンテーションの精度向上を目的とする。",
  "application_field": "医療画像解析",
  "problem_evidence": "low contrast and irregular tumor boundaries",
  "solution_evidence": "modified U-Net architecture",
  "confidence": "high"
}
\end{verbatim}

\section{確認方針}
LLMによる抽出結果は，そのまま正解として扱わない．A1，B1，および補助ケースから代表サンプルを選び，problem\_contextが課題のみを表しているか，solution\_contextが解決手段のみを表しているかを人手で確認する．確認は，○，△，×などの判定で行う．本研究では，抽出精度を完全に保証するのではなく，研究上利用できる程度の妥当性があるかを確認する．
